\documentclass[10pt,letterpaper]{article}
% Add a bunch of useful math, font, and symbols
\usepackage{amsfonts}
\usepackage{amsmath}
\usepackage{amssymb}

% English support
\usepackage[english]{babel}

% Citation
\usepackage[superscript]{cite}

% Better enumerate and itemize
\usepackage{enumitem}

% Better control of headers and footers
\usepackage{fancyhdr}

% Floating objects like figures and tables
\usepackage{float}

% Page layout and dimensioning
\usepackage[margin=1in]{geometry}

% Basic color, graphics and text manipulation 
\usepackage{graphicx}

% Use Helvetica typeface
\usepackage[scaled]{helvet}
\renewcommand\familydefault{\sfdefault} 
\usepackage[T1]{fontenc}

% Cross-referencing hyperlinks
\usepackage{hyperref}

% Line break for long URLs
\usepackage{breakurl}

% Accept utf-8 input encoding
\usepackage[utf8]{inputenc}

% Make indexes
\usepackage{makeidx}

% Microtype (apparently makes the typographics stuff better)
\usepackage{microtype}

% [Disabled] multi-column writing
% \usepackage{multicol}

% English ordinal counting (1st, 2nd, etc.)
\usepackage{nth}

% Long table
\usepackage{longtable}

% Paragraph skip - adds extra lineskip spacing
\usepackage{parskip}
\setlength{\parskip}{0.7\baselineskip plus 2pt}

% Add ability to set space between lines
\usepackage{setspace}

% S.I. units
\usepackage{siunitx}

% Subcaptions for subfigures
\usepackage{subcaption}

% Include svg graphics
\usepackage{svg}

% Drawing graphics
\usepackage{tikz}

% Subsubsubsection
\usepackage{titlesec}
\setcounter{secnumdepth}{4}
\titleformat{\paragraph}
{\normalfont\normalsize\bfseries}{\theparagraph}{1em}{}
\titlespacing*{\paragraph}
{0pt}{3.25ex plus 1ex minus .2ex}{1.5ex plus .2ex}

% Custom titles
\usepackage{titling}

% Url
\usepackage{url}

\RequirePackage[figure,table]{totalcount}


% Custom definitions
\newcommand{\doctitle}{Validation Specification}
\newcommand{\docsubtitle}{Computing Platform Multirotor with FPGA Hardware Acceleration Applications}

% Custom commands
\newcommand{\ts}{\textsubscript}	% Subscript command %


% Use hyphans to break up urls
\def\UrlBreaks{\do\/\do-}

% PDF and href setup
% Hyper ref
\hypersetup{
	colorlinks=true,
	citecolor=black,
	linkcolor=black,
	filecolor=black,
	urlcolor=blue,
	pdftitle={\@title},
	bookmarks=true
}
\urlstyle{same}
% Page headings
\pagestyle{fancy}
\fancyhead[L]{\MakeUppercase{CPEN/ELEC 491}}
\fancyhead[C]{\textbf{\doctitle}}
\fancyhead[R]{Mieszko Lis, PhD}
\fancyfoot{}
\fancyfoot[C]{Non-Confidential}
\fancyfoot[R]{\thepage}

% No paragraph indent
\parindent 0ex

% Meta
\author{
	Deutsch, Peter &
	\textit{me@peterdeutsch.ca}
	\\
	He, Muchen &
	\textit{m.he@alumni.ubc.ca}
	\\
	Hsueh, Arthur &
	\textit{ah11962@outlook.com}
	\\
	Wang, Meng &
	\textit{wzfftxwd@gmail.com}
	\\
	Wilson, Ardell &
	\textit{ardellw96@gmail.com}
}
\title{\doctitle}
\date{\today}
\makeatletter
\renewcommand{\maketitle}{
	\bgroup
	\setlength{\parindent}{0pt}
	\begin{flushleft}
		% Top spacing
		\vspace*{0.5in}

		% Team logo
		\includegraphics[scale=0.4]{../assets/capstonelogo1.png}
		\vspace*{0.25in}

		% Title
		\textbf{\Huge{\@title}}\\
		\hrulefill

		% Subtitle
		\textbf{\huge{\docsubtitle}}
		
		\vspace*{0.5in}

		% Course number and team
		\textbf{\Large{CPEN/ELEC 491 Capstone Team 109}}\\
		\hspace*{0.1cm}
		\begin{tabular}[h]{|ll}
			\@author
		\end{tabular}

		\vspace*{0.25in}

		\textbf{\Large{Mieszko Lis, PhD}}\\
		\hspace*{0.1cm}
		\begin{tabular}[h]{|ll}
			Electrical and Computer Engineering, The University of British Columbia
		\end{tabular}

		\vfill
		
		% Date
		\large{Revision 1.0 -- \@date}
		\vspace*{0.5in}

		% Logo
		\hspace*{-0.3cm}\includegraphics[scale=0.5]{../assets/ece_logo.pdf}

	\end{flushleft}
	\egroup
}
\makeatother

% Begin Document
\begin{document}

% Title Page
\begin{titlepage}
	\maketitle
\end{titlepage}

% Executive Summary (Not used here in proposal)
% \include{executive_summary.tex}

% Revision history
%\thispagestyle{empty}
\section*{Revision History}
The full revision history and commited changes of the document can be found in the git repository history: \href{https://github.com/Capstone-Skynet/Capstone-Skynet.github.io}{https://github.com/Capstone-Skynet/Capstone-Skynet.github.io/commits/master}.

\begin{table}[H]
\begin{tabular}{*{4}{l}p{0.5\linewidth}}
\hline
Version \# & Initials & Release Date & Changeset & Changes Made \\ \hline

0.0 & PD & 2019-10-11 & \texttt{660e001} & Initial skeleton of the document.\\
0.1 & MH & 2019-10-11 & \texttt{6af9e8a} & Populate initial document with draft content required for Milestone I.\\
0.2 & PD & 2019-11-23 & & Initial framework for test descriptions created.\\
0.3 & MW & 2019-11-24 & & First set of tests added.\\
1.0 & PD & 2019-11-25 & & General clean-up and release for Milestone II.\\
1.1 & AH & 2020-02-09 & & Added Full System Tests \\

 & & & \\ \hline
\end{tabular}
\end{table}

% Table of contents
\setcounter{secnumdepth}{3}
\tableofcontents

% Terms and Abbreviations
\thispagestyle{empty}

\section*{Terms and Abbreviations}

Technical terms and abbreviations dictionary go here.

\begin{tabular}[h]{rp{0.75\linewidth}}
    \hline
    \textbf{Term} & \textbf{Definition}\\
    \hline

    ANN & Artificial Neural Network, or simply ``Neural Network'', is a data processing model modeled after neuron interactions. The process consists of forward propagation using several matrix multiplications.\cite{ann}\\
    ASIC & Application-specific Integrated Circuits.\\
    CNN & Convolutional neural network are neural networks that is especially useful for image classification.\cite{cnn} \\
    ECE & Department of Electrical and Computer Engineering at the University of British Columbia.\\
    FPGA & Field-Programmable-Gate-Arrays, ``programmable'' hardware that allows ASIC-like performance with software-like turn-around time and flexibility.\\
    GPU & Graphics Processing Unit, a discrete piece of hardware designed to accelerate graphic-intensive or other parallel computing tasks.\\
    LOS & Line-of-sight.\\
    ML & Machine learning.\\
    Multirotor & An unmanned vehicle with multiple engines. \\
    OTS & Off-the-shelf, or commercially available/purchasable \\
    PID / PID Controller & Proportion-Integral-Derivative controllers is the most common control algorithm for precise and accurate movement, as well as to compensate external forces.\cite{pid}\\
    RNN & Recurrent neural networks are neural networks where the output depends on previous computations, essentially consists of memory.\cite{rnn}\\
    RX & Receiver.\\
    TC & Transport Canada.\\
    TX & Transmitter.\\
    YOLO & You-Only-Look-Once is a fast ML algorithm that detect objects but is unlike CNN nor RNN.\cite{yolo}\cite{yolo-2}\\
     & \\

    \hline

\end{tabular}


% List of figures and tables
\thispagestyle{empty}
\listoffigures
\listoftables
\newpage

% Set page and section counter
\setcounter{page}{1}


% TODO: fill out all the sections
% if the sections gets too long, move them to a separte .tex document
\section{About This Document}

\subsection{Purpose}
This document serves to outline the methods, tests, and scenarios which are used to validate the project, and the status of these validation steps. 

\subsection{Intended Audience}
This document is intended for:
\begin{itemize}
\item \textbf{Quality Controllers (the Capstone team)}: to assess the validation methodology and results,
\item \textbf{Designers (the Capstone team)}: to assess the current degree of completion of the project,
\item \textbf{Client (and their representatives)}: to assess the final state of the project (when complete),
\item \textbf{and Legal Personnel:} to review the validation scheme if necessary. Such personnel may include personnel from the Department of Electrical and Computer Engineering, Industry Canada, and Transport Canada
\end{itemize}

\subsection{Reading Guide}
This document organizes tests based on their component sub-systems, and whether the tests are automated or not (see Section \ref{auto} for a detailed description). Each section has a table which details the test, including a unique identifier, the requirement which it is testing, the test's status, and the last time the test was executed.

\section{Validation Overview}

\subsection{Automated and Manual Tests}\label{auto}
The validation tests employed (and this document) are split into two categories, \textit{automated tests} and \textit{manual tests}:

\textbf{Automated tests} are tests focused on tasks which are performed automatically and regularly by the device and \textit{can be monitored by the device itself}. Examples of such scenarios include self-testing systems on boot, monitoring communication status, machine learning accuracy logging, et cetera.

\textbf{Manual tests} are tests which \textit{cannot be performed by the device itself} and must be manually performed by a quality controller. These tests are performed less frequently --- some may be executed before every flight, while others may only be performed once. 

\subsection{Test Notation}
The following notations and conventions are used throughout test tables in this document:

\textbf{Verification Requirement (VR)}: The unique identifier of a specific test. These identifiers always start with \textbf{T}.

\textbf{Design Requirement (DR)}: Denotes the design requirement which the test is exercising, as listed in the project's \textit{Requirements Specification}. Requirement identifiers are encoded by their first characters, as follows:
\begin{itemize}
	\item \textbf{F}: Functional Requirement,
	\item \textbf{NF}: Non-Functional Requirement, and
	\item \textbf{C}: Constraint
\end{itemize}

\textbf{Status}: The status of a specific test:
\begin{itemize}
	\item \textbf{PEND} if the specific test has never been performed,
	\item \textbf{FAIL} if the specific test is failing, or
	\item \textbf{PASS} if the specific test is passing
\end{itemize}

\textbf{Effective Date}: Denotes the date at which the test was last performed (and thus, when the status was last updated). This field is not applicable for pending tests.

\textbf{It is important to note that a one-to-one mapping of VRs to DRs \textit{does not necessarily exist}}. Some requirements (particularly those which are high-level or functional) will have no tests, while some requirements will have multiple tests. It is also possible for one test to exercise multiple requirements at once.

\section{Automated Test Cases}

\subsection{Computing Platform Tests}

\subsection{Multirotor Tests}

\section{Manual Test Cases}

\subsection{Computing Platform Tests}
\subsubsection{General Tests}

\subsubsection{Camera Tests}
\subsubsection{Machine Learning Tests}
\subsubsection{Base Station Tests}
\subsubsection{Power Tests}

\subsection{Multirotor Tests}
\paragraph{T.DR.1: RPAS Maximum Takeoff Weight Test}

This test obtains the maximum thrust output, minimum thrust-to-weight ratio, and maximum power output of the RPAS prototype.

\textbf{Addresses}:  NF.DR.1, NF.DR.12, NF.DR.13

\textbf{Procedure}:
\begin{enumerate}[noitemsep]
    \item Attach the RPAS prototype to a scale or force sensor and ensure that the RPAS cannot separate scale or force sensor.
    \item Attach multimeter to the RPAS power circuitry.
    \item Fix or constrain the RPAS to the ground by using a heavy weight, tape, or fasteners.
    \item Tare the scale or force-sensor such that the measure is reference to the initial conditions.
    \item Power on the drone and arm the motors.
    \item Apply maximum throttle from the radio transmitter.
    \item Obtain voltage and current measurements from the multimeter.
    \item Obtain force measurements from the scale or force sensor.
\end{enumerate}

\textbf{Verification}: 
To ensure that the RPAS has minimum manuverability, at maximum take-off weight, we should be using 80\% throttle. For this test to pass, the RPAS must achieve a maximum pull of 125\% (2.5 kg) of maximum takeoff weight. Since this is a maximum load test case, we verify the measured current and voltage to ensure NF.DR.12 and NF.DR.13 are checked.

% 

\paragraph{T.DR.2: RPAS Typical Takeoff Weight Test}

The typical takeoff weight test quantifies the flight-time of the RPAS during a typical operation with significantly less massive payload.

\textbf{Addresses}: NF.DR.7, NF.DR.10

\textbf{Procedure}:
\begin{enumerate}[noitemsep]
    \item Attach voltage and current monitor to the RPAS power circuitry.
    \item Set up RPAS for test-flight in open area or large room.
    \item Power on the drone and arm the motors.
    \item Start timer.
    \item Apply reasonable throttle from the radio transmitter such that the RPAS is hovering 3.0 m above ground.
    \item Record the \%-throttle applied for hovering.
    \item Keep hovering until the RPAS battery reaches 10\% capacity.
    \item Stop timer when the RPAS lands; calculate flight time in minutes.
    \item Measure the total mass of the RPAS and payload, and compute the average power consumption using the data acquired from the voltage and current monitors. Calculate efficiency in g/W.
\end{enumerate}

\textbf{Verification}:
The measured flight time should be greater than 20 minutes with a payload mass of 500 g in order to satisfy the requirement NF.DR.7. The calculated efficiency must be greater than 5.0 g/W to satisfy the requirement NF.DR.10.

\section{Verification Results}
\subsection{Multirotor Verification Results}
\begin{table}[H]
	\centering
	\begin{tabular}{llll}
	\hline
	\textbf{VR} & \textbf{DR} & \textbf{Status} & \textbf{Eff. Date} \\
	\hline
	T.DR.1 & NF.DR.1, NF.DR.12, NF.DR.13  & PEND   & \\
	T.DR.2 & NF.DR.7, NF.DR.10 & PEND   & \\
	\hline
	\end{tabular}
\end{table}

% Bibliography
\clearpage
\addcontentsline{toc}{section}{References}
\bibliographystyle{ieeetr}
\bibliography{references}

% Appendix (uncomment to enable appendix)
% \clearpage
% \appendix
% \section{Appendix name}\label{appendix:sample-appendix}
% Content here

\end{document}
